\begin{frame}{Atividade final}
  \begin{block}{Desafio: Git Wall}
    Vamos colocar em prática o que aprendemos fazendo uma
    contribuição ao projeto \textbf{GitWall}.
  \end{block}
  \begin{columns}[T]
    \begin{column}{0.48\textwidth}
      \begin{block}{O que fazer}
        \begin{enumerate}
          \item Faça um \textit{fork} do projeto.
          \item Faça sua contribuição seguindo as instruções do arquivo \texttt{README.md}.
          \item Registre as alterações com Git e envie para o GitHub.
          \item Abra um \textit{Pull Request}.
        \end{enumerate}
      \end{block}
    \end{column}
    \begin{column}{0.48\textwidth}
      \begin{block}{Entrega}
        O \textit{Pull Request} é a entrega da atividade
        para obtenção do certificado.
        \vspace{0.3cm} \\
        \centering
        \href{https://github.com/silash35/git-wall}
        {\textbf{🌐 github.com/silash35/git-wall}}
      \end{block}
    \end{column}
  \end{columns}
  \vspace{0.3cm}
  \begin{block}{}
    Não quer aparecer no GitWall?
    Basta sinalizar isso no seu \textit{Pull Request}.
  \end{block}
\end{frame}

\begin{frame}{Material extra}
  \begin{block}{Documentação oficial}
    \href{https://git-scm.com/learn}{\textbf{🌐 git-scm.com/learn}}
    \vspace{0.2cm} \\
    Tutoriais, vídeos e o livro \textit{Pro Git}.
  \end{block}
  \begin{block}{Git Cheat Sheet}
    \href{https://git-scm.com/cheat-sheet}{\textbf{🌐 git-scm.com/cheat-sheet}}
    \vspace{0.2cm} \\
    Resumo dos principais comandos e operações do Git.
  \end{block}
  \begin{block}{Curso de Git e GitHub}
    \href{https://www.youtube.com/playlist?list=PLx4x_zx8csUhopU45en_VDUS9hkze64xO}
    {\textbf{▶ Curso de Git e GitHub — YouTube}}
    \vspace{0.2cm} \\
    Curso de Git GitHub do CFBCursos
  \end{block}
\end{frame}

\begin{frame}{Pra explorar depois}
  \begin{columns}[T]
    \begin{column}{0.5\textwidth}
      \begin{itemize}
        \item Git Rebase
        \item Git Stash
        \item Cherry-pick
        \item Git Revert vs Reset
        \item Git Flow
        \item Git LFS
        \item Git Submodules
        \item Squash de commits
      \end{itemize}
    \end{column}
    \begin{column}{0.5\textwidth}
      \begin{itemize}
        \item Git Reflog
        \item Tags e Releases
        \item GitHub Actions
        \item GitHub Pages
        \item Issues e Projects
        \item Git Hooks
        \item Conventional Commits
      \end{itemize}
    \end{column}
  \end{columns}
\end{frame}